% Einheit 2 von 2 (Katalog-Einheit 5): Diagramme beurteilen und Boxplot
\einheitenkopf[e2][2 Diagramme beurteilen und Boxplot]{Einheit 2 von 2 · Diagramme beurteilen und Boxplot}
\zweigzeile{Hier lernst du, Aussagen zu Diagrammen zu prüfen, Täuschungen zu erkennen und Boxplots zu lesen · neu oder schon bekannt – je nach Buch (am Gymnasium neu in diesem Jahr) · P10 · baut auf: Säulendiagramme lesen (Nr. 4), Kenngrößen (Nr. 8), Bruchteile (Nr. 9), Prozentsatz (Nr. 1)}
\setcounter{aufgabe}{24}

\begin{aufgabe}{Ich erkenne, wo die Achse anfängt.}
Kreuze an, ob die senkrechte Achse bei 0 beginnt. Wenn nicht, umkreise den Startwert der Achse.

\noindent\begin{minipage}[t]{0.48\linewidth}
\begin{teile}\teil \kreuz{beginnt bei 0}\quad\kreuz{beginnt nicht bei 0}\end{teile}
\saeulenab[ymax=40,ystep=10,yfein=10,ylabel=Anzahl]{A/30,B/20}
\end{minipage}\hfill
\begin{minipage}[t]{0.48\linewidth}
\begin{teile}\teil \kreuz{beginnt bei 0}\quad\kreuz{beginnt nicht bei 0}\end{teile}
\saeulenab[ymin=60,ymax=80,ystep=5,yfein=5,ylabel=Preis in €]{2023/65,2024/75}
\end{minipage}

\noindent\begin{minipage}[t]{0.48\linewidth}
\begin{teile}\teil \kreuz{beginnt bei 0}\quad\kreuz{beginnt nicht bei 0}\end{teile}
\saeulenab[ymin=1000,ymax=1400,ystep=100,yfein=100,ylabel=Mitglieder]{2023/1150,2024/1300}
\end{minipage}\hfill
\begin{minipage}[t]{0.48\linewidth}
\begin{teile}\teil \kreuz{beginnt bei 0}\quad\kreuz{beginnt nicht bei 0}\end{teile}
\saeulenab[ymax=200,ystep=50,yfein=50,ylabel=Besucher]{Sa/150,So/100}
\end{minipage}
\end{aufgabe}

\verfahren{Aussagen zu einem Diagramm prüfen}
\begin{aufgabe}{Ich kann eine Aussage zu einem Diagramm prüfen.}
Ein Zoo zählt jedes Jahr seine Besucher.

\liniendia[ymax=60,ystep=10,yfein=2,ylabel=Besucher in Tausend]{2019/56,2020/50,2021/44,2022/46,2023/40,2024/36}

Stimmt die Aussage?
\begin{teile}
\teil 2019 kamen 56 Tausend Besucher. \hfill \janein
\teil 2023 kamen 44 Tausend Besucher. \hfill \janein
\teil 2021 kamen weniger Besucher als 2022. \hfill \janein
\teil 2020 kamen 50 Besucher. \hfill \janein
\teil Von 2023 bis 2024 ging die Zahl um 4 Tausend zurück. \hfill \janein
\teil Die Zahl der Besucher ist von 2019 bis 2024 jedes Jahr gesunken. (P10 2022 OS) \hfill \janein
\end{teile}
\end{aufgabe}

\begin{aufgabe}{Ich kann Aussagen mit Anteilen zu einem Diagramm prüfen.}
180 Schülerinnen und Schüler der 7. Klassen nennen ihr Lieblingsfach, jede und jeder genau eines.

\balkenab[xmax=80,xstep=10,xfein=2,xlabel=Nennungen]{Sport/76,Mathematik/40,Kunst/30,Deutsch/22,Musik/12}

Kreuze an und begründe mit einer Rechnung.
\begin{teile}
\teil Mehr als die Hälfte nennt Sport.\\ \kreuz{wahr}\quad\kreuz{falsch}\quad\kreuz{nicht entscheidbar}
\teil Mathematik wird um ein Drittel öfter genannt als Kunst.\\ \kreuz{wahr}\quad\kreuz{falsch}\quad\kreuz{nicht entscheidbar}
\teil Jeder 15. nennt Musik.\\ \kreuz{wahr}\quad\kreuz{falsch}\quad\kreuz{nicht entscheidbar}
\teil 15\,\% nennen Musik.\\ \kreuz{wahr}\quad\kreuz{falsch}\quad\kreuz{nicht entscheidbar}
\teil Mehr Mädchen als Jungen nennen Sport.\\ \kreuz{wahr}\quad\kreuz{falsch}\quad\kreuz{nicht entscheidbar}
\teil Sport wird am häufigsten genannt, und Deutsch wird halb so oft genannt wie Mathematik. (P10 2018 OS)\\ \kreuz{wahr}\quad\kreuz{falsch}\quad\kreuz{nicht entscheidbar}
\end{teile}
\end{aufgabe}

\verfahren{Verzerrte Diagramme erkennen}
\begin{aufgabe}{Ich kann erklären, warum eine abgeschnittene Achse täuscht.}
Ein Sportverein zeigt seine Mitgliederzahlen im linken Diagramm.

\noindent\begin{minipage}[t]{0.48\linewidth}
\saeulenab[ymin=400,ymax=460,ystep=10,yfein=5,ylabel=Mitglieder]{2023/420,2024/440}
\end{minipage}\hfill
\begin{minipage}[t]{0.48\linewidth}
\saeulenab[ymax=500,ystep=100,yfein=20,ylabel=Mitglieder]{2023/,2024/}
\end{minipage}

\begin{teile}
\teil Die Säule für 2024 ist doppelt so hoch wie die für 2023. Hat sich die Zahl der Mitglieder verdoppelt? Begründe.
\teil Um wie viel Prozent ist die Zahl der Mitglieder gestiegen? Runde auf eine Stelle. \hfill \leerfeld[\%]
\teil Zeichne die beiden Säulen in das rechte Diagramm, dessen Achse bei 0 beginnt. Vergleiche mit dem linken.
\end{teile}
\end{aufgabe}

\begin{aufgabe}{Ich kann erklären, warum eine gedehnte Achse täuscht.}
Beide Diagramme zeigen den Preis für eine Kugel Eis.

\noindent\begin{minipage}[t]{0.48\linewidth}
Diagramm A\par
\liniendia[ymax=6,ystep=1,yfein=0.5,ylabel=Preis in €,breite=0.9]{2020/1.2,2021/1.3,2022/1.4,2023/1.5,2024/1.6}
\end{minipage}\hfill
\begin{minipage}[t]{0.48\linewidth}
Diagramm B\par
\liniendia[ymax=1.6,ystep=0.2,yfein=0.1,ylabel=Preis in €,breite=0.9]{2020/1.2,2021/1.3,2022/1.4,2023/1.5,2024/1.6}
\end{minipage}

\begin{teile}
\teil Lies in Diagramm B den Preis für 2020 und für 2024 ab. \hfill \leerfeld[€] \quad \leerfeld[€]
\teil In welchem Diagramm scheint der Preis stärker zu steigen? Woran liegt das?
\teil Um wie viel Prozent ist der Preis von 2020 bis 2024 gestiegen? Runde auf eine Stelle. \hfill \leerfeld[\%]
\end{teile}
\end{aufgabe}

\begin{aufgabe}{Ich kann die Fortschreibung eines Trends prüfen. (kommt in Klasse 9)}
Ein Laden zählt die verkauften DVDs.

\liniendia[ymax=600,ystep=100,yfein=20,ylabel=verkaufte DVDs]{2019/520,2020/440,2021/360,2022/280,2023/200,2024/160}

\begin{teile}
\teil Beschreibe den Trend.
\teil Der Besitzer sagt: „Wenn das so weitergeht, verkaufen wir 2026 keine einzige DVD mehr.“ Prüfe die Aussage.
\end{teile}
\end{aufgabe}

\begin{aufgabe}{Ich kann die Fortschreibung eines Trends prüfen – weiter.}
Eine Zeitung schreibt: „Kinobesuche brechen ein – jedes Jahr gehen weniger Menschen ins Kino.“ Dazu zeigt sie dieses Diagramm (Besuche in einem Land).

\saeulenab[ymin=90,ymax=130,ystep=10,yfein=2,ylabel=Besuche in Mio.]{2019/122,2020/112,2021/116,2022/104,2023/98,2024/96}

\begin{teile}
\teil Prüfe die Aussage der Zeitung und erkläre, warum das Diagramm den Rückgang stärker zeigt, als er ist. (P10 2019 OS)
\end{teile}
\end{aufgabe}

\verfahren{Boxplot}
\begin{aufgabe}{Ich kann einen Boxplot lesen. (kommt in Klasse 8, je nach Buch bis Kl. 9)}
Die 7a hat beim Weitsprung gemessen.

\begin{boxplots}[xmin=240,xmax=480,xstep=10,xlabel=Weite in cm]
\bp{260,310,340,380,450}{7a}
\end{boxplots}

Lies ab.
\begin{teilezwei}
\tz kürzeste Weite: \leerfeld[cm] & \tz größte Weite: \leerfeld[cm] \\
\tz Median: \leerfeld[cm] & \\
\end{teilezwei}
\begin{teile}
\teil unteres und oberes Quartil: \leerfeld[cm] \quad \leerfeld[cm]
\end{teile}

\anweisung{Bestimme.}
\begin{teile}
\teil Spannweite: \leerfeld[cm]
\teil Wie viel Prozent der Kinder sprangen mindestens 310 cm weit? \hfill \leerfeld[\%]
\end{teile}
\end{aufgabe}

\begin{aufgabe}{Ich kann einen Boxplot zeichnen.}
Zeichne den Boxplot in die Zeile mit demselben Buchstaben.
\begin{teile}
\teil Minimum 3, unteres Quartil 6, Median 9, oberes Quartil 12, Maximum 18
\teil Punkte, schon geordnet: 4, 6, 7, 9, 11, 12, 14, 17
\teil Punkte: 15, 8, 11, 5, 19, 10, 13, 7, 12, 16, 9, 14
\end{teile}

\begin{boxplots}[xmin=0,xmax=20,xstep=1,xlabel=Punkte]
\bp{}{c)}
\bp{}{b)}
\bp{}{a)}
\end{boxplots}
\end{aufgabe}

\begin{aufgabe}{Ich kann zwei Boxplots vergleichen.}
Die Boxplots zeigen die Punkte beim Mathetest in zwei Klassen.

\begin{boxplots}[xmin=0,xmax=20,xstep=1,xlabel=Punkte]
\bp{9,11,12,14,16}{7b}
\bp{6,10,13,15,19}{7a}
\end{boxplots}

\begin{teile}
\teil Welche Klasse hat den höheren Median? \hfill \leerfeld
\teil In welcher Klasse ist die Spannweite größer? \hfill \leerfeld
\teil In welcher Klasse liegt die mittlere Hälfte der Ergebnisse enger beieinander? \hfill \leerfeld
\teil Tom sagt: „Die 7a war besser, denn ihr bestes Ergebnis ist höher.“ Beurteile die Aussage mit zwei Kenngrößen.
\end{teile}
\end{aufgabe}

\begin{aufgabe}{Ich finde den Fehler bei der Deutung eines Diagramms.}
Finde den Fehler, benenne ihn und korrigiere.
\begin{teile}
\teil Die Achse eines Säulendiagramms beginnt bei 80. Säule A zeigt 85 verkaufte Eis, Säule B 95 verkaufte Eis. Ben sagt: „Säule B ist dreimal so hoch wie Säule A, also wurden bei B dreimal so viele Eis verkauft.“
\par\vspace{16mm}
\teil Ein Boxplot hat das Minimum 12, das untere Quartil 20, den Median 26, das obere Quartil 35 und das Maximum 41. Mia sagt: „Die Box geht von 20 bis 35, also ist die Spannweite 15.“
\par\vspace{16mm}
\end{teile}
\end{aufgabe}

\begin{aufgabe}{Ich kann begründen, warum die Achse eines Säulendiagramms bei 0 beginnen sollte.}
Begründe.
\begin{teile}
\teil Warum sollte die senkrechte Achse eines Säulendiagramms bei 0 beginnen?
\par\vspace{12mm}
\teil In einem Liniendiagramm zur Körpertemperatur eines kranken Kindes beginnt die Achse bei 36\,°C. Ist das eine Täuschung?
\par\vspace{12mm}
\end{teile}
\end{aufgabe}

\begin{aufgabe}{Ich kann Diagramme in der Werbung beurteilen.}
Zwei Handyanbieter zeigen ihre Kundenzahlen.

\noindent\begin{minipage}[t]{0.48\linewidth}
Anbieter A\par
\saeulenab[ymin=950,ymax=1000,ystep=10,yfein=5,ylabel=Kunden in Tausend]{2023/960,2024/990}
\end{minipage}\hfill
\begin{minipage}[t]{0.48\linewidth}
Anbieter B\par
\saeulenab[ymax=500,ystep=100,yfein=20,ylabel=Kunden in Tausend]{2023/400,2024/440}
\end{minipage}

\begin{teile}
\teil Lies die Kundenzahlen ab und berechne für beide Anbieter den Zuwachs von 2023 bis 2024 in Tausend und in Prozent. Runde auf eine Stelle.
\teil Anbieter A wirbt: „Unsere Kundenzahl ist explodiert!“ Beurteile die Werbung.
\end{teile}
\end{aufgabe}
